\section{Diskussion und Reflexion}
\label{sec:diskussion}

Die Diskussion und Reflexion der Ergebnisse zeigt, dass die Prognose des Aluminiumpreises eine herausfordernde Aufgabe darstellt, die durch hohe Volatilität und die Vielzahl von Einflussfaktoren geprägt ist. Die Messungen deuten darauf hin, dass das entwickelte System zwar in der Lage ist, gewisse Muster zu erkennen, jedoch die Vorhersagbarkeit über den 1-Tages-Horizont hinaus stark abnimmt.

\paragraph{Markteffizienz und Vorhersagbarkeit}
Das geringe Bestimmtheitsmaß ($R^2 \approx 0,0043$) und eine Richtungsgenauigkeit von $51,1 \%$ stehen im Einklang mit der \emph{Efficient Market Hypothesis} (EMH). Nach dieser Theorie spiegeln Marktpreise alle öffentlich verfügbaren Informationen wider, sodass Preisänderungen auf neue, unvorhersehbare Informationen reagieren und somit der Wahrscheinlichkeitstheorie (Random Walk) folgen. Die Ergebnisse deuten darauf hin, dass die historischen Preis- und Volumendaten allein nicht ausreichen, um systematische Vorhersagemuster zu extrahieren. Dies könnte auf die hohe Komplexität des Marktes zurückzuführen sein, in dem zahlreiche Faktoren (z. B. geopolitische Ereignisse, technologische Entwicklungen) eine Rolle spielen, die nicht direkt in den Zeitreihen erfasst werden.

\paragraph{Bedeutung der probabilistischen Prognose}
Ein Aspekt des Systems ist die Abkehr von algorithmischen Punktvorhersagen hin zu probabilistischen Verteilungen. Mittels MC-Dropout wird die Modellunsicherheit erkannt. Die Coverage von $90,8 \%$ für den 1-Tages-Horizont zeigt, dass das Modell in der Lage ist, die Varianz der Daten in einem statistischen Rahmen abzubilden. Dieser Ansatz bietet eine Grundlage für das Risikomanagement, da er keine exakte Punktlandung verspricht, sondern ein Spektrum der wahrscheinlichen Preisentwicklungen definiert. In Phasen mit hoher Unsicherheit (weite Konfidenzintervalle) kann die Beschaffungsstrategie  angepasst werden, um Risiken zu minimieren.

\paragraph{Grenzen der Datenbasis}
Die Beschränkung auf historische Preis- und Volumenzeitreihen sowie gesamtwirtschaftliche Indizes lässt weitere Faktoren außer Acht. Variablen wie physische Lagerbestände (z. B. LME Stocks), Energiepreise als Primärkostenfaktor der Produktion sowie politische Entscheidungen fließen nur indirekt über die Preisbildung in das Modell ein. Eine Erweiterung der Datenbasis um diese Zeitreihen oder unstrukturierte Daten (Text-Mining von Marktberichten) bietet Potenzial für eine Steigerung der Signalqualität.
